% !PRINTED 414
\weberpage{414}
%<<CONT
can be so chosen that the constant $\varepsilon$ occurring in (10) is \gesperrt{always}
a third root of unity. If $\pi$ belongs to the $\alpha$, then $\varepsilon$
is at the same time a second root of unity and must therefore be $=1$.
From this we infer the theorem:

2. \gesperrt{The function $\vartheta$ remains unchanged under the
permutations of the four-group $A$.}

The third power of $\vartheta$ admits the permutations of the alternating
group, and $\vartheta$ is therefore the root of a pure cubic equation
$\vartheta^3 = z$. This, then, is the resolvent $\Phi(\vartheta,z) = 0$,
which in this case is a \gesperrt{partial resolvent} (Bd. I, \S.\,161).

From these considerations there now results without difficulty the following theorem,
whose proof is the goal of this whole
investigation:

3. \gesperrt{A general equation of degree higher than
4 has no resolvents with one parameter.}

For if $n$ is greater than 4, then, if we regard $\vartheta$ as a function
of any four of the variables $x$, we can apply Theorem 2.
It then follows that $\vartheta$ remains unchanged when any
pair of transpositions is performed among the variables. Since one
can now compose the whole alternating group out of transposition pairs
(Bd. I, \S.\,153, 8.), it follows that $\vartheta$ remains unchanged by
the alternating group altogether, and is therefore an alternating or
a symmetric function. $\Phi(\vartheta,z) = 0$ reduces to
the identity $\vartheta = z$ and is not a resolvent.

Let it further be remarked that, if the symmetric group is put in place
of the alternating group, this theorem holds a fortiori.

The theorem whose proof we have developed here was first
stated by \gesperrt{Kronecker}, without his publishing a proof
of it. The first proof was given by \gesperrt{F. Klein}, which
was then further simplified by \gesperrt{Gordan}\footnote{\gesperrt{F. Klein}, Vorlesungen über das Ikosaëder. The
proof by \gesperrt{Gordan}, which our presentation follows in its main
features, is found in Bd. 29 der mathematischen Annalen (1887).}.
